\section{Differential Equations}

\subsection{Modelling and Definitions}

\begin{remark}
    \textbf{Modelling Process}: To derive an ODE from a physical scenario:
    \begin{enumerate}
        \item Identify the quantity $u(t)$ and express its change over a small interval $\Delta t$:
        $u(t+\Delta t) = u(t) + (\text{rate}) \cdot \Delta t$
        \item Rearrange and take the limit $\Delta t \to 0$:
        $\frac{u(t+\Delta t) - u(t)}{\Delta t} = \text{rate} \implies u'(t) = \text{rate}$
    \end{enumerate}
\end{remark}

\begin{remark}
    \textbf{Typical ODE Models:}
    \begin{enumerate}
        \item Bounded Growth (Logistic): With capacity $L$:
        $$u'(t) = k u(t) \left(1 - \frac{u(t)}{L}\right)$$
        \item Bimolecular Reaction ($A + B \to C$): With concentrations $a, b$:
        $$u'(t) = k(a - u(t))(b - u(t))$$
    \end{enumerate}
\end{remark}

\begin{definition}
    An \textbf{ordinary differential equation (ODE)} contains derivatives of a function of one independent variable. It is:
    \begin{itemize}
        \item \textbf{of $n$-th order}: if the $n$-th derivative is the highest occurring.
        \item \textbf{linear}: if the unknown function $u(t)$ and its derivatives appear only as a linear combination.
        \item \textbf{homogeneous}: if it only contains terms with $u(t)$ and its derivatives (otherwise \textbf{inhomogeneous}).
    \end{itemize}
\end{definition}

\begin{theorem}
    \textbf{Picard-Lindelöf (Existence and Uniqueness)}: If $f(x, y)$ is continuous and Lipschitz continuous with respect to $y$, then the IVP $y' = f(x, y)$ with $y(x_0) = y_0$ has a unique solution in some interval around $x_0$.
\end{theorem}

\subsection{First-Order ODEs}

\begin{theorem}
    \textbf{Separation of Variables:} For $y' = f(x)g(y) \implies \int \frac{dy}{g(y)} = \int f(x)dx + C$.
\end{theorem}

\begin{theorem}
    \textbf{Linear First-Order ODE:} $y' + p(x)y = q(x)$.
    \textbf{Integrating Factor Method:} 
    1. Calculate $F(x) = \exp\left(\int p(x)dx\right)$.
    2. Multiply ODE by $F(x)$: $(y F(x))' = q(x)F(x)$.
    3. Integrate: $y(x) = \frac{1}{F(x)} \left(\int q(x)F(x)dx + C\right)$.
\end{theorem}

\begin{theorem}
    \textbf{Bernoulli ODE:} $y' + p(x)y = q(x)y^n$ ($n \neq 0, 1$).
    \textbf{Trick:} Substitute $u = y^{1-n} \implies u' = (1-n)y^{-n}y'$.
    This transforms the ODE into the linear form: $\frac{u'}{1-n} + p(x)u = q(x)$.
\end{theorem}

\begin{theorem}
    \textbf{Homogeneous substitution:} $y' = f(y/x)$. Substitute $u = y/x \implies y = ux, y' = u'x + u$.
\end{theorem}

\subsection{Linear ODEs of n-th Order}

\begin{theorem}
    \textbf{Superposition Principle}: If $y_1$ and $y_2$ solve a linear homogeneous ODE, then for any $C_1, C_2 \in \mathbb{R}$:
    $$y(x) = C_1 y_1(x) + C_2 y_2(x)$$
\end{theorem}

\begin{theorem}
    An $n$-th order linear homogeneous ODE has $n$ linearly independent solutions $y_1, \dots, y_n$. Its general solution is:
    $$y(x) = \sum_{i=1}^n C_i y_i(x) \quad \text{with } C_i \in \mathbb{R}$$
\end{theorem}

\begin{lemma}
    For a homogeneous ODE with constant coefficients $a_n y^{(n)} + \dots + a_0 y = 0$, the Ansatz $y(x) = e^{\lambda x}$ yields the characteristic equation:
    $$a_n \lambda^n + a_{n-1} \lambda^{n-1} + \dots + a_1 \lambda + a_0 = 0$$
    with characteristic polynomial $p(\lambda)$.
\end{lemma}

\begin{theorem}
    \textbf{General Solution from Characteristic Roots}: The roots of $p(\lambda)$ determine the fundamental solutions:
    \begin{itemize}
        \item \textbf{Distinct real root $\lambda \in \mathbb{R}$}
        $$y(x) = e^{\lambda x}$$
        \item \textbf{Multiple real root $\lambda \in \mathbb{R}$ with multiplicity $m$}
        $$e^{\lambda x}, x e^{\lambda x}, \dots, x^{m-1} e^{\lambda x}$$
        \item \textbf{Complex pair $\lambda = \alpha \pm i\beta \in \mathbb{C}$ (multiplicity $m$)}
        $$e^{\alpha x}\cos(\beta x), e^{\alpha x}\sin(\beta x), \dots,$$
        $$x^{m-1}e^{\alpha x}\cos(\beta x), x^{m-1}e^{\alpha x}\sin(\beta x)$$
    \end{itemize}
\end{theorem}

\subsection{Inhomogeneous Linear ODEs}

\begin{lemma}
    The general solution of $a_n y^{(n)} + \dots + a_0 y = s(x)$ is $y(x) = y_h(x) + y_p(x)$, where $y_h$ is the general homogeneous solution and $y_p$ is a particular solution.
\end{lemma}

\subsubsection{Method of Undetermined Coefficients (Ansatz Table)}
For $ay'' + by' + cy = s(x)$, if no resonance:
\begin{center}
\scriptsize
\setlength{\tabcolsep}{2pt}
\begin{tabular}{|l|l|}
\hline
\textbf{Source $s(x)$} & \textbf{Ansatz $y_p(x)$} \\ \hline
$P_n(x)$ (Polyn.) & $Q_n(x) = $ \\
& $A_n x^n + \dots + A_0$ \\ \hline
$A e^{kx}$ & $B e^{kx}$ \\ \hline
$A \sin(\omega x)$ & $C \sin(\omega x)$ \\
$+ B \cos(\omega x)$ & $+ D \cos(\omega x)$ \\ \hline
$e^{\alpha x} P_n(x)$ & $e^{\alpha x} Q_n(x)$ \\ \hline
$e^{\alpha x} \sin(\beta x)$ & $e^{\alpha x} (C \cos(\beta x)$ \\
& $+ D \sin(\beta x))$ \\ \hline
\end{tabular}
\end{center}
\textbf{Resonance Rule:} If the standard Ansatz $y_p(x)$ contains a solution to the homogeneous equation, multiply the \textit{entire} Ansatz by $x$ (or $x^2$ for double roots).

\begin{important}
    \textbf{Resonance}: If the term $e^{kx}$ or $e^{i\omega x}$ in $s(x)$ matches an $m$-fold root of the characteristic polynomial $p(\lambda)$, multiply the standard Ansatz $y_p(x)$ by $x^m$.
\end{important}

\subsubsection{Variation of Constants}

\begin{theorem}
    \textbf{Variation of Constants (First-Order Linear ODE)}: For $y' + p(x)y = q(x)$, the general solution is:
    $$y(x) = e^{-P(x)} \left( \int q(x) e^{P(x)} \, dx + C \right)$$
    where $P(x) = \int p(x) \, dx$ is an antiderivative of $p(x)$.
\end{theorem}

\subsection{Initial Value Problems and Boundary-Value Problems}

\begin{definition}
    An \textbf{Initial Value Problem (IVP)} specifies conditions at a single point $x_0$ (e.g. $y(x_0) = y_0$, $y'(x_0) = y_1$).
    A \textbf{Boundary Value Problem (BVP)} specifies conditions at multiple points (e.g. $y(a) = y_a$, $y(b) = y_b$).
\end{definition}
\subsection{Exam Strategy: ODE Solver Workflow}
\begin{enumerate}
    \item \textbf{Can I separate?} $y' = f(x)g(y)$.
    \item \textbf{Is it linear 1st order?} $y' + p(x)y = q(x) \to$ Integrating factor $e^{\int p}$.
    \item \textbf{Is it 2nd order linear with constant coeffs?} $ay'' + by' + cy = s(x)$.
        \begin{itemize}
            \item Solve char. eq. $a\lambda^2 + b\lambda + c = 0$.
            \item Particular solution: Match $s(x)$ with Ansatz table (check for resonance!).
        \end{itemize}
\end{enumerate}

\subsection{Exam Strategy: Constrained Extrema}
To find extrema of $f(x,y)$ s.t. $g(x,y)=0$:
\begin{enumerate}
    \item Set up system $\nabla f = \lambda \nabla g$ and $g = 0$.
    \item Solve for $x, y, \lambda$.
    \item \textbf{Don't forget the boundary!} If the region is a disk $x^2+y^2 \le 1$, check interior $\nabla f = 0$ AND boundary $x^2+y^2=1$.
\end{enumerate}
